% Drafted introduction and related work — to be rewritten in unified voice by writing agent

\section{Introduction}

Single-cell RNA sequencing (scRNA-seq) has transformed cancer biology by exposing the transcriptional heterogeneity of malignant cells and their tumour microenvironment at unprecedented resolution~\cite{tang2009mrnaseq,patel2014single,tirosh2016dissecting}. Beyond gene expression, scRNA-seq signals also carry a coarse-grained record of the underlying genomic landscape: large copy number amplifications and deletions (CNAs) shift average transcript levels across affected chromosomal intervals, making them detectable from read depth alone~\cite{patel2014single,beroukhim2010landscape}. Because CNAs are central drivers of tumour initiation, progression, and therapy resistance~\cite{beroukhim2010landscape,stephens2011chromothripsis,mcgranahan2017clonal}, inferring them directly from scRNA-seq offers a powerful route to study clonal structure in the same cells whose phenotypes are being measured.

Single-cell DNA sequencing (scDNA-seq) remains the gold standard for CNA profiling at single-cell resolution~\cite{navin2011tumor,gawad2016singlecell}, but its cost and sparse coverage limit routine use. Several computational methods now aim to recover CNA profiles from scRNA-seq alone. These methods fall into two methodological families. The first relies solely on gene-expression signals: inferCNV performs a corrected moving average over gene windows~\cite{tickle2019infercnv}; CopyKAT couples integrative Bayesian segmentation with hierarchical clustering to identify normal reference cells and delineate clonal substructure~\cite{gao2021copykat}; and SCEVAN applies multichannel segmentation based on the Mumford--Shah energy principle after selecting high-confidence normal cells via tumour microenvironment gene signatures~\cite{defalco2023scevan}. The second family augments expression with allelic information. Numbat uses a haplotype-aware hidden Markov model that fuses expression, allelic ratios from scRNA-seq reads, and population haplotype phasing to call allele-specific CNAs~\cite{gao2023numbat}, while CaSpER employs a multiscale signal-processing framework that integrates B-allele frequency (BAF) shifts with smoothed expression profiles~\cite{serin2020casper}.

Despite the proliferation of these tools, their comparative behaviour under realistic conditions has not been systematically characterised. Existing evaluations are largely self-reported by each method's original authors and rely on heterogeneous ground-truth sources. Independent, head-to-head benchmarks against paired scDNA-seq from the \emph{same} cell have been missing, leaving practitioners without clear guidance on which tool to choose for their dataset, how tumour purity and sequencing depth shape accuracy, and whether these tools detect loss of heterozygosity (LOH) reliably.

Recent advances in single-cell multi-omics~\cite{macaulay2015gtseq,zachariadis2020dntrseq,zhou2020singlecell} now make such direct comparison possible. Joint measurement of DNA and RNA in individual cells provides a ground-truth CNA profile for every cell whose transcriptome is being analysed. We exploit this opportunity to carry out a comprehensive, independent benchmark of the five most widely used scRNA-seq-based CNA inference tools.

Our study evaluates inferCNV~\cite{tickle2019infercnv}, CopyKAT~\cite{gao2021copykat}, SCEVAN~\cite{defalco2023scevan}, Numbat~\cite{gao2023numbat}, and CaSpER~\cite{serin2020casper} across thirteen samples spanning colorectal cancer (CRC), glioma, acute lymphoblastic leukaemia (ALL), neuroendocrine tumour, a gastric sample, and NPC43 and HUVEC cell lines~\cite{zhou2020singlecell,zachariadis2020dntrseq}. We assess four canonical tasks --- tumour versus normal cell classification, CNA profile accuracy, subclonal structure inference, and aneuploidy detection in non-malignant cells --- and systematically perturb tumour purity (T:N ratios from 1:100 to 100:1) and sequencing depth (median 10K, 3K, 1K UMIs per cell). We also examine the role of reference-cell selection, including the effect of including tumour microenvironment (TME) cells, and evaluate LOH detection against allele-specific copy number states derived from paired scDNA-seq using Sequenza~\cite{favero2015sequenza} and DNAcopy-based circular binary segmentation~\cite{olshen2004dnacopy,willenbrock2005comparison}. Our contributions are: (i) an independent head-to-head evaluation against paired scDNA-seq ground truth; (ii) quantitative characterisation of how tumour purity, sequencing depth, and reference selection reshape each tool's accuracy; (iii) a detailed audit of Numbat's LOH calls revealing that $\sim$82\% correspond to non-neutral CNV events; and (iv) practical recommendations for tool selection and parameter tuning, including pairings that maximise robustness.
